%%%%%%%%%%%%%%%%%%%%%%%%%%%%%%%%%%%%%%%%%%%%%%%%%%%%%%%%%%%
%  Resumen de Enfoque Radiológico — Tórax PA
%  Basado en tau-class v2.6.0 (CC BY 4.0)
%%%%%%%%%%%%%%%%%%%%%%%%%%%%%%%%%%%%%%%%%%%%%%%%%%%%%%%%%%%



\documentclass[9pt,a4paper,twocolumn,twoside]{tau-class/tau}
\usepackage[spanish,es-nodecimaldot,es-noindentfirst]{babel}

%----------------------------------------------------------
% Datos del documento
%----------------------------------------------------------

\doctype{Técnicas Radiológicas}
\title{Tórax Frente PA}

%----------------------------------------------------------
% Header / Footer
%----------------------------------------------------------

\footinfo{Tórax}
\theday{\today}
\organization{Tórax}
\leadauthor{Técnicas Radiológicas}

%----------------------------------------------------------

\begin{document}
	
	\maketitle
	\thispagestyle{firststyle}
	
	%==========================================================
	\section{Tórax}
	%==========================================================
	
	\subsection{Tórax Frente PA}
	
	\subsubsection*{1. Identificación}
	
	\textbf{Tipo:} Proyección básica estándar. Forma par con la Proyección Lateral de Tórax.\\
	
	\textbf{Indicaciones clínicas principales:}
	\begin{itemize}
		\item Neumonía, bronquitis e infecciones del tracto respiratorio inferior
		\item Neumotórax (se complementa con proyección en espiración forzada)
		\item Evaluación de silueta cardiaca: cardiomegalia, derrame pericárdico
		\item Derrame pleural y engrosamiento pleural
		\item Lesiones mediastínicas y masas pulmonares
		\item Control post-quirúrgico, post-intubación y seguimiento de patologías crónicas
	\end{itemize}
	
	%------------------------------------------------------
	\subsubsection*{2. Técnica Base}
	
	\textbf{Receptor de imagen:} 35 × 43 cm. Orientación \textbf{transversal} en paciente hiperesténico y esténico; \textbf{longitudinal} en asténico e hipoesténico. Borde superior 4 a 5 cm por encima de los hombros.
	
	\textbf{DFRI:} 180 a 200 cm. Minimiza la magnificación de estructuras mediastínicas por divergencia del haz y optimiza la radioprotección.
	
	\textbf{Técnica:} Ultra homogeneizada (5 o 6 homogenizaciones). Produce escala de grises larga para visualizar simultáneamente trama pulmonar, mediastino y estructuras óseas, minimizando contraste innecesario y dosis al paciente.
	
	\textbf{Factores técnicos:} kVp \textbf{110} — mAs \textbf{2,5}. Espesor medido a los \textbf{2/3 del tórax} (plano de los pechos). Fórmula base: mAs = 10 × 3,24 × 3,5 $\approx$ 114, dividido por 2 en cada homogenización. La pérdida de mAs se compensa con Bucky y mayor DFRI.
	
	\textbf{Bucky:} Bucky mural (vertical) con grilla antidifusora.
	
	%------------------------------------------------------
	\subsubsection*{3. Posicionamiento}
	
	\paragraph{Paciente}
	Bipedestación o sedestación. Plano medio sagital perpendicular al receptor y centrado en su línea media. Siempre en PA: el corazón y el mediastino quedan cerca del receptor reduciendo su magnificación; protege tiroides, tejido mamario y esternón (órgano hematopoyético primario). Pies bien apoyados y separados. El ortostatismo desciende el diafragma por gravedad, evita la ingurgitación vascular pulmonar y permite demostrar niveles hidroaéreos.
	
	\paragraph{Región de interés}
	Hombros descendidos y en contacto firme con el receptor, ambos en el mismo plano transversal, lo que desciende las clavículas y despeja los vértices. Brazos en jarra: dorso de las manos sobre las caderas, codos orientados hacia el receptor; esto rota internamente los húmeros y desplaza las escápulas fuera de los campos pulmonares. Mentón extendido levemente hacia arriba para no bloquear vías aéreas ni vértices. En pacientes con mamas voluminosas, pedirle que las eleve y separe lateralmente para despejar las bases pulmonares.
	
	\paragraph{Rayo Central}
	Perpendicular al plano del receptor (0°). Punto de entrada: \textbf{T7}, usando el ángulo inferior de las escápulas como reparo anatómico.
	
	\paragraph{Receptor de imagen}
	Centrado en T7. Borde superior 4 a 5 cm por encima de los hombros para incluir los vértices. El plano medio sagital del paciente coincide con la línea media del receptor.
	
	\paragraph{Respiración}
	Inspiración forzada máxima, sostenida durante la exposición. Garantiza la máxima expansión del tórax en los tres ejes y el descenso del diafragma. En sospecha de neumotórax: dos exposiciones, una al final de inspiración forzada y otra al final de espiración forzada. La comparación evidencia aire libre pleural oculto en inspiración; también demuestra movilidad diafragmática, cuerpo extraño y atelectasia.
	
	%------------------------------------------------------
	\subsubsection*{4. Anatomía Visible}
	
	Tráquea aireada en línea media. Ambos campos pulmonares completos desde vértices hasta ángulos costofrénico y cardiofrénico. Cúpulas diafragmáticas. Silueta cardiaca y botón aórtico. Trama vascular pulmonar y árbol bronquial. Costillas, clavículas, escápulas proyectadas fuera de los campos y columna torácica. La trama pulmonar debe ser visible a través de la silueta cardiaca. Los ángulos costofrénico y cardiofrénico deben visualizarse íntegros.
	
	%------------------------------------------------------
	\subsubsection*{5. Criterios de Evaluación}
	
	\textbf{Posicionamiento correcto:} campos pulmonares íntegros incluyendo vértices y ángulos costofrénico y cardiofrénico; \textbf{9 a 10 arcos costales posteriores} (horizontales) y \textbf{6 a 7 arcos costales anteriores} (curvos) visibles por encima de las cúpulas; contornos definidos del corazón y el diafragma; marcas pulmonares desde el hilio a la periferia; marca de lado derecho en la imagen.
	
	\textbf{Ausencia de rotación:} extremidades mediales de ambas clavículas equidistantes de la columna (si la izquierda parece más corta, el paciente está rotado hacia la izquierda); apófisis espinosas centradas en la línea media; tráquea en la línea media.
	
	\textbf{Penetración y contraste:} trama pulmonar visible a través de la silueta cardiaca, confirmando la técnica ultra homogeneizada. Pulmones, mediastino y estructuras óseas diferenciados en la misma imagen.
	
	%------------------------------------------------------
	\subsubsection*{6. Puntos Críticos}
	
	\textbf{Rotación del paciente:} asimetría de clavículas y distorsión de la sombra cardiaca. Corrección: pedir que adelante el hombro más alejado del Bucky y lo presione hasta que ambos queden en contacto firme y parejo. Si una clavícula aparece más alta, indicarle que al inspirar no eleve los hombros y los relaje en un plano horizontal paralelo al piso.
	
	\textbf{Escápulas superpuestas:} brazos no en jarra o codos sin orientarse hacia adelante. La rotación interna del húmero es la que desplaza las escápulas.
	
	\textbf{Inspiración insuficiente:} menos de 9 costillas posteriores visibles, campos condensados y corazón aparentemente agrandado. Indicador de inspiración correcta: 9--10 costillas posteriores y 6--7 anteriores, con trama pulmonar visible a través del corazón.
	
	\textbf{Corte de ángulos costofrénico o cardiofrénico:} pérdida de información diagnóstica crítica; repetir siempre que sea posible.
	
	\textbf{Adaptaciones:} paciente que no puede estar de pie: AP en sedestación o decúbito, la mayor magnificación cardiaca es esperada y debe consignarse. Mamas voluminosas: elevarlas y separarlas. Clavícula congénitamente más alta: no confundir con rotación.
	
	\textbf{Trucos:} espesor medido a los 2/3 del tórax, no en el punto más grueso. La DFRI larga compensa la pérdida de mAs de las homogenizaciones. Si no se ve trama a través del corazón, el kVp o las homogenizaciones fueron insuficientes. La PA magnifica menos el corazón que la AP, relevante para medir el índice cardiotorácico.
	
	%----------------------------------------------------------
	
\end{document}